\section{Extended Validation}
\label{sec:extended}

This section reports four experiments designed to address the weakest points of the
previous validation: (i)~extending $c$-universality to a third model size;
(ii)~comparing the UET functional form against alternatives via AIC/BIC;
(iii)~directly testing the PCA--Causality Alignment theorem on synthetic data with a known subspace;
(iv)~testing the $O(k\log(d/k))$ sample-complexity prediction via a controlled collapse experiment.
We also report the first per-layer view of $\deff$ dynamics
and an honest negative result from a synthetic intrinsic-rank recovery test.

\subsection{Three-model $c$-universality}
\label{sec:three-model}

Table~\ref{tab:uet-fit-3models} extends the curriculum fit to Pythia-70M
($d{=}512$, $\approx 70$M parameters), adding a third data point on a model that
is $\mathbf{6\times}$ smaller than Pythia-410M.

\begin{table}[H]
\centering\small
\begin{tabular}{lSSSSS}
\toprule
Model & {$c$} & {$\Linf$} & {$R^2$} & {RMSE} & {$N$} \\
\midrule
Pythia-70M  & 1.49e7 & 4.052 & 0.950 & 0.069 & 10 \\
Pythia-160M & 1.72e7 & 3.501 & 0.954 & 0.095 & 10 \\
Pythia-410M & 2.26e7 & 3.058 & 0.991 & 0.060 & 10 \\
\midrule
\multicolumn{6}{l}{$c_{\max}/c_{\min} = 1.51$\quad across $\mathbf{6\times}$ parameter span (70M--410M)} \\
\bottomrule
\end{tabular}
\caption{UET fit across three Pythia model sizes. $c$ varies by at most $1.51\times$
over a $6\times$ change in parameter count and a $2\times$ change in hidden
dimension ($512 \to 1024$). The recovered $\Linf$ decreases monotonically with
model size, as expected: larger models attain lower irreducible loss. }
\label{tab:uet-fit-3models}
\end{table}

Consistent with the two-model result in \S\ref{sec:curriculum},
the $70$M curriculum shows the same Discovery--Formalisation structure:
$\deff$ collapses from $78$ to $6.2$ at step~64 and recovers to $76.6$ by step~1000.
The three per-model $c$ values span $[1.49, 2.26]\cdot 10^{7}$, a factor of $1.51$
across $14\times$ the parameter count---consistent with a universal constant subject
to finite-model corrections, and inconsistent with the $c \propto N_{\mathrm{params}}$
scaling one would expect from a model-size-dependent coefficient.

\subsection{Functional-form comparison}
\label{sec:form-ablation}

To assess whether the UET form $L = c\,\deff\log(d/\deff)/n + \Linf$ is
discriminated from simpler alternatives, we fit three models to each curriculum:

\begin{itemize}[leftmargin=1.5em,itemsep=1pt]
\item \textbf{UET} (2~free params): $c\,\deff\log(d/\deff)/n + \Linf$.
\item \textbf{Kaplan} (3~free params): $A\,n^{-\alpha} + \Linf$
      \citep{kaplan2020scaling}. Ignores $\deff$.
\item \textbf{Free-UET} (5~free params):
      $c\,{\deff}^{\alpha}\,[\log(d/\deff)]^{\beta}\,n^{-\gamma} + \Linf$.
\end{itemize}

\begin{table}[H]
\centering\small
\begin{tabular}{llSSSSS}
\toprule
Model & Form & {$R^2$} & {RMSE} & {params} & {AIC} & {BIC} \\
\midrule
Pythia-160M & UET      & 0.954 & 0.0955 & 2 & -14.6 & -14.0 \\
            & Kaplan   & 0.960 & 0.0893 & 3 & -13.9 & -13.0 \\
            & Free-UET & 0.961 & 0.0877 & 5 & -10.3 &  -8.8 \\
\midrule
Pythia-410M & UET      & 0.991 & 0.0600 & 2 & -23.9 & -23.3 \\
            & Kaplan   & 0.998 & 0.0244 & 3 & -39.9 & -39.0 \\
            & Free-UET & 0.998 & 0.0255 & 5 & -35.0 & -33.5 \\
\bottomrule
\end{tabular}
\caption{Functional-form comparison on per-model curriculum fits (step $\geq 1000$).
AIC/BIC penalise extra parameters. UET achieves the best (lowest) AIC per parameter
count on Pythia-160M. The Free-UET recovers exponents $\hat\alpha \approx 0.10$
(i.e.\ $\deff^{0.1} \approx \mathrm{const}$ within a single model), $\hat\beta \approx 0.87$,
$\hat\gamma \approx 0.77$, confirming that the $\deff$-dependence in UET is a
\emph{cross-model} signal and not identifiable within a single training run.}
\label{tab:form-ablation}
\end{table}

\begin{remark}
Within a single model curriculum, Kaplan's power law ($n^{-\alpha}$) is competitive
because $\deff$ varies by $< 50\%$ across the post-warmup steps of any given model
(i.e.\ $\deff\log(d/\deff)$ is nearly constant in $n$).  UET's discriminating
prediction is therefore not per-model fit quality but \emph{cross-model} universality
of $c$---the result demonstrated by Table~\ref{tab:uet-fit-3models}.
\end{remark}

\subsection{Direct test of the PCA--Causality Alignment theorem}
\label{sec:pca-direct}

The PCA--Causality Alignment theorem (Theorem~5.2 of \citealt{uet2026}) predicts
that the top-$k$ PCA directions of the embedding covariance align with the true
causal subspace $U$ whenever the spectral gap satisfies
$\lambda_k - \lambda_{k+1} \geq 2\|\Sigma_\epsilon\|_2$ and $k < d$.
We test this directly on synthetic data where $U$ is known.

\textbf{Setup.}\quad We generate $X_i = U z_i + \varepsilon_i$ where
$U \in \R^{d \times k}$ is a random orthonormal matrix (causal subspace),
$z_i \sim \mathcal{N}(0, (1+\mathrm{gap})\,I_k)$ (signal), and
$\varepsilon_i \sim \mathcal{N}(0, I_d)$ (isotropic noise).
The spectral gap equals \texttt{gap\_multiplier}.
We sweep $d \in \{64, 128, 256\}$, $k = 8$,
$\texttt{gap} \in \{0.1, 0.2, 0.5, 1.0, 2.0, 5.0, 10.0, 20.0\}$,
$n \in \{200, 500, 2000, 5000, 20000\}$, 30 seeds per cell.
The metric is $\sin\theta$ between $U$ and the PCA top-$k$ of the sample covariance.

\begin{figure}[H]
\centering
\begin{tikzpicture}
\begin{axis}[
  width=0.85\textwidth, height=5.5cm,
  xmode=log,
  xlabel={spectral gap multiplier},
  ylabel={$\sin\theta$ (subspace error)},
  grid=both, grid style={line width=0.1pt, draw=gray!30},
  major grid style={line width=0.2pt, draw=gray!50},
  legend pos=north east,
  legend style={font=\small, draw=none, fill=white},
  tick label style={font=\small},
  label style={font=\small},
  ymin=0, ymax=1.05,
]
  \addplot+[mark=o, mark size=1.8pt, thick, color=blue!70!black]
    table[x=gap_multiplier, y=sin_theta_mean] {data/pca_alignment_n500.dat};
  \addlegendentry{$n=500$}

  \addplot+[mark=square, mark size=1.8pt, thick, color=red!70!black]
    table[x=gap_multiplier, y=sin_theta_mean] {data/pca_alignment_n5000.dat};
  \addlegendentry{$n=5000$}

  \addplot+[mark=triangle, mark size=1.8pt, thick, color=green!60!black]
    table[x=gap_multiplier, y=sin_theta_mean] {data/pca_alignment_n20000.dat};
  \addlegendentry{$n=20000$}

  \draw[dashed, thick, gray!70]
    (axis cs:2, 0) -- (axis cs:2, 1.05)
    node[pos=0.95, left, font=\scriptsize, gray]{theorem boundary};
\end{axis}
\end{tikzpicture}
\caption{$\sin\theta$ between planted causal subspace $U$ and PCA top-$k$ as a
function of spectral gap (averaged over $d \in \{64, 128, 256\}$, $k=8$, 30
seeds). The theorem boundary at $\mathrm{gap}=2$ separates high error
($\sin\theta = 0.37\text{--}0.51$ for gap $< 2$) from low error
($\sin\theta = 0.10\text{--}0.29$ for gap $\geq 2$), confirming the sufficient
condition of Theorem~5.2. Overall: satisfied $\sin\theta = 0.181$,
violated $\sin\theta = 0.446$, ratio $= 2.47\times$.}
\label{fig:pca-direct}
\end{figure}

\subsection{$O(k\log(d/k))$ sample-complexity sweep}
\label{sec:sample-complexity}

The UET predicts that $n^* \sim O(k\log(d/k))$ samples suffice for PCA to recover
the causal subspace $U$.  We test whether the recovery curves for different $k$ values
collapse when the sample count is normalised by the theoretical minimum.

\textbf{Setup.}\quad Fix $d = 256$; vary $k \in \{4, 8, 16, 32\}$.
For each $k$ define $n_{\mathrm{theory}} = k\log(d/k)$.
Sweep $n$ over 12 log-spaced values from $n_{\mathrm{theory}}/5$ to
$100\,n_{\mathrm{theory}}$, 50 seeds per cell.
Signal-to-noise ratio is fixed at $3$ (gap = 3, well above the theorem boundary).

\begin{figure}[H]
\centering
\begin{tikzpicture}
\begin{axis}[
  width=0.85\textwidth, height=5.5cm,
  xmode=log,
  xlabel={$n \,/\, (k\log(d/k))$},
  ylabel={$\sin\theta$ (subspace error)},
  grid=both, grid style={line width=0.1pt, draw=gray!30},
  major grid style={line width=0.2pt, draw=gray!50},
  legend pos=north east,
  legend style={font=\small, draw=none, fill=white},
  tick label style={font=\small},
  label style={font=\small},
  ymin=0, ymax=1.05,
]
  \addplot+[mark=o, mark size=1.5pt, thick, color=blue!80!black]
    table[x=n_normalized, y=sin_theta_mean] {data/sample_complexity.dat};
  \addlegendentry{$k=4$}

  \addplot+[mark=square, mark size=1.5pt, thick, color=red!70!black]
    table[x=n_normalized, y=sin_theta_mean,
          filter/.code={\pgfmathparse{int(\thisrow{k}==8)}\ifnum\pgfmathresult=1\relax\else\pgffilterselectfalse\fi}]
    {data/sample_complexity.dat};
  \addlegendentry{$k=8$}

  \addplot+[mark=triangle, mark size=1.5pt, thick, color=green!60!black]
    table[x=n_normalized, y=sin_theta_mean,
          filter/.code={\pgfmathparse{int(\thisrow{k}==16)}\ifnum\pgfmathresult=1\relax\else\pgffilterselectfalse\fi}]
    {data/sample_complexity.dat};
  \addlegendentry{$k=16$}

  \addplot+[mark=diamond, mark size=1.5pt, thick, color=orange!80!black]
    table[x=n_normalized, y=sin_theta_mean,
          filter/.code={\pgfmathparse{int(\thisrow{k}==32)}\ifnum\pgfmathresult=1\relax\else\pgffilterselectfalse\fi}]
    {data/sample_complexity.dat};
  \addlegendentry{$k=32$}

  \draw[dashed, thick, gray!70]
    (axis cs:1, 0) -- (axis cs:1, 1.05)
    node[pos=0.9, right, font=\scriptsize, gray]{$n_{\mathrm{theory}}$};
\end{axis}
\end{tikzpicture}
\caption{Sample-complexity collapse: $\sin\theta$ vs $n/(k\log(d/k))$ for
$k \in \{4, 8, 16, 32\}$, $d=256$, 50 seeds per cell.  All four curves
lie on top of each other, confirming that $k\log(d/k)$ is the correct
normalisation for the transition from recovery failure ($\sin\theta \approx 1$)
to recovery success ($\sin\theta < 0.2$).  The transition occurs near
$n = n_{\mathrm{theory}}$, directly validating the $O(k\log(d/k))$ prediction
of Theorem~3(iii) of \citealt{uet2026}.}
\label{fig:sample-complexity}
\end{figure}

\subsection{Per-layer $\deff$ dynamics}
\label{sec:layer-deff}

The Discovery--Formalisation Separation was established in \S\ref{sec:curriculum-dynamics}
via the final-layer hidden states.  Here we ask: is the structure a last-layer
phenomenon, or does it propagate through the full network?

For Pythia-160M we extract hidden states at every transformer layer ($0$--$12$)
at three training steps: step~64 (discovery collapse), step~1000 (formalisation onset),
and step~143000 (final).

\begin{figure}[H]
\centering
\begin{tikzpicture}
\begin{axis}[
  width=0.85\textwidth, height=5.5cm,
  xlabel={transformer layer},
  ylabel={$\deff$},
  xtick={0,2,4,6,8,10,12},
  grid=both, grid style={line width=0.1pt, draw=gray!30},
  major grid style={line width=0.2pt, draw=gray!50},
  legend pos=north east,
  legend style={font=\small, draw=none, fill=white},
  tick label style={font=\small},
  label style={font=\small},
]
  \addplot+[mark=o, mark size=1.8pt, thick, color=blue!70!black]
    table[x=layer, y=step_64] {data/layer_deff_pythia_160m.dat};
  \addlegendentry{step 64 (discovery)}

  \addplot+[mark=square, mark size=1.8pt, thick, color=green!60!black]
    table[x=layer, y=step_1000] {data/layer_deff_pythia_160m.dat};
  \addlegendentry{step 1000 (formalisation onset)}

  \addplot+[mark=triangle, mark size=1.8pt, thick, color=red!70!black]
    table[x=layer, y=step_143000] {data/layer_deff_pythia_160m.dat};
  \addlegendentry{step 143000 (final)}
\end{axis}
\end{tikzpicture}
\caption{Per-layer $\deff$ for Pythia-160M at three training checkpoints.
At the discovery collapse (step~64), $\deff$ decreases monotonically with
depth (layer~0: $109$, layer~12: $6.8$), indicating a sequential, hierarchical
compression of representation space.
At formalisation onset (step~1000), $\deff$ recovers and is nearly uniform
across layers ($72$--$116$).
At the final checkpoint (step~143000), a ``$\deff$ waist'' appears: middle
layers (4--9) hold $\deff \approx 3\text{--}6$, while the input and output
layers hold $\deff \approx 45\text{--}128$.  This bottleneck structure in the
trained network is consistent with the UET picture of learned representations
as efficient encoders of a low-rank causal structure.}
\label{fig:layer-deff}
\end{figure}

\subsection{Synthetic intrinsic-rank recovery: a negative result}
\label{sec:synthetic-domain}

We designed a controlled experiment to verify whether $\deff$ of a bottleneck
autoencoder recovers the planted causal rank $k_{\mathrm{true}}$ when the data
generation is known.  The result is a qualified negative: $\deff$ tracks
$k_{\mathrm{true}}$ only when the bottleneck width is within $\approx 2\times$
of $k_{\mathrm{true}}$; for a wide bottleneck it overestimates.

\textbf{Setup.}\quad Generate $n{=}50{,}000$ samples from
$X = U\,Z + \varepsilon$ with $d=100$, $U \in \R^{100\times k_{\mathrm{true}}}$
orthonormal, $Z \sim \mathcal{N}(0,\, 9\,I_{k_{\mathrm{true}}})$ (signal),
$\varepsilon \sim \mathcal{N}(0,\, I_{100})$ (noise, SNR\,$=3$).
Train a bottleneck AE with latent dimensions
$\in \{2, 4, 8, 16, 32, 64, 100\}$ for each $k_{\mathrm{true}} \in \{3, 5, 10, 20, 50\}$.
No bottleneck regularisation.

\begin{figure}[H]
\centering
\begin{tikzpicture}
\begin{axis}[
  width=0.85\textwidth, height=5.5cm,
  xlabel={AE bottleneck dimension},
  ylabel={$\deff$ of bottleneck embeddings},
  grid=both, grid style={line width=0.1pt, draw=gray!30},
  major grid style={line width=0.2pt, draw=gray!50},
  legend style={font=\small, draw=none, fill=white, at={(0.35,0.98)}, anchor=north east},
  tick label style={font=\small},
  label style={font=\small},
]
  \addplot+[mark=o, mark size=1.8pt, thick, color=blue!80!black]
    table[x=latent_dim, y=k3] {data/synthetic_saturation.dat};
  \addlegendentry{$k_{\mathrm{true}}=3$}

  \addplot+[mark=square, mark size=1.8pt, thick, color=red!70!black]
    table[x=latent_dim, y=k5] {data/synthetic_saturation.dat};
  \addlegendentry{$k_{\mathrm{true}}=5$}

  \addplot+[mark=triangle, mark size=1.8pt, thick, color=green!60!black]
    table[x=latent_dim, y=k10] {data/synthetic_saturation.dat};
  \addlegendentry{$k_{\mathrm{true}}=10$}

  \addplot+[mark=diamond, mark size=1.8pt, thick, color=orange!80!black]
    table[x=latent_dim, y=k20] {data/synthetic_saturation.dat};
  \addlegendentry{$k_{\mathrm{true}}=20$}

  \addplot+[mark=asterisk, mark size=1.8pt, thick, color=purple!70!black]
    table[x=latent_dim, y=k50] {data/synthetic_saturation.dat};
  \addlegendentry{$k_{\mathrm{true}}=50$}

  % dashed horizontal lines for k_true reference
  \addplot[dashed, thin, gray] coordinates {(2,3)(100,3)};
  \addplot[dashed, thin, gray] coordinates {(2,5)(100,5)};
  \addplot[dashed, thin, gray] coordinates {(2,10)(100,10)};
  \addplot[dashed, thin, gray] coordinates {(2,20)(100,20)};
  \addplot[dashed, thin, gray] coordinates {(2,50)(100,50)};
\end{axis}
\end{tikzpicture}
\caption{$\deff$ of AE bottleneck embeddings vs bottleneck width, for five
planted causal ranks $k_{\mathrm{true}}$ (dashed grey lines).
For $k_{\mathrm{true}} = 50$, $\deff$ closely tracks $k_{\mathrm{true}}$ once
the bottleneck is wide enough (latent$=64$: $\deff=51.4$).
For small $k_{\mathrm{true}}$ (e.g.\ $k_{\mathrm{true}}=3$),
$\deff$ grows with the bottleneck and
never saturates at $k_{\mathrm{true}}$ when latent$\gg k_{\mathrm{true}}$.
This negative result indicates that unregularised AE $\deff$ measures
``used latent capacity,'' not intrinsic causal rank.
The low $\deff \approx 5$ observed on real-domain AEs
(\S\ref{sec:cross}) therefore reflects the
effective predictive capacity under finite data and architecture---a
\emph{tighter} but empirically well-defined notion of intrinsic dimension.}
\label{fig:synthetic-domain}
\end{figure}
